\subsubsection{Variational Principle}\label{sec:variational-principle}

This section presents a \textbf{theorem that makes VQE possible}. Recall the logical chain established in previous sections: we need to find the ground state of a Hamiltonian, which is equivalent to determining its minimum eigenvalue. Now, the problem is: ``\emph{\textbf{how can we find the minimum eigenvalue of a Hamiltonian?}}''

\highspace
\textcolor{Red2}{\faIcon{exclamation-triangle} \textbf{Diagonalization is not feasible for large Hamiltonians!}} An $n$-qubit quantum system lives in a Hilbert space of dimension $2^{n}$, therefore the Hamiltonian matrix has size $2^{n} \times 2^{n}$. For example, a system of $n=100$ qubits would require a Hamiltonian matrix of size $2^{100} \times 2^{100}$, which is approximately $10^{30} \times 10^{30}$. One naïve approach would be to find the minimum eigenvalue of the Hamiltonian by \textbf{diagonalizing} it. Diagonalizing a matrix involves \textbf{computing all its eigenvalues and eigenvectors}. However, the \hl{computational cost grows polynomially with the matrix size}, while the \hl{matrix size itself grows exponentially with the number of qubits.} Therefore, the \textbf{overall computation becomes exponential}.

\highspace
\begin{flushleft}
    \textcolor{Green3}{\faIcon{check-circle} \textbf{Variational Principle saves us}}
\end{flushleft}
Wait a minute, do we really need $\lambda_0$, $\lambda_1$, $\lambda_2$, ..., $\lambda_{2^n-1}$? No, we only need the \textbf{minimum eigenvalue} $\lambda_0$. So, instead of solving the complete eigenvalue problem (finding all eigenvalues), we only want to find the \textbf{minimum eigenvalue}.

\highspace
The \definitionWithSpecificIndex{Rayleigh-Ritz Variational Principle}{Rayleigh-Ritz Variational Principle}{Variational Principle}\footnote{%
    The word ``variational'' comes from the fact that we \textbf{vary} the trial state $\ket{\psi}$ to minimize the expectation value $\bra{\psi} H \ket{\psi}$. We keep varying the parameters until the energy is as low as possible.%
} saves us from diagonalization. It states that:
\begin{equation*}
    \bra{\psi} H \ket{\psi} \geq \lambda_0 \qquad \text{for any normalized state } \ket{\psi}
\end{equation*}
We remark that the expectation value $\bra{\psi} H \ket{\psi}$ (\autopageref{sec:expectation-value}) represents the \textbf{average value obtained if we repeatedly measure the observable $H$ on many identical copies of the state $\ket{\psi}$}. In other words, it is the expected energy of the quantum state $\ket{\psi}$.

\highspace
\begin{theorem}[Rayleigh-Ritz Variational Principle]
    Let $H$ be a Hermitian operator (Hamiltonian) with eigenvalues $\lambda_0 \leq \lambda_1 \leq \dots \leq \lambda_{2^n-1}$ and corresponding eigenvectors $\ket{\psi_0}, \ket{\psi_1}, \dots, \ket{\psi_{2^n-1}}$. For \textbf{any normalized state} $\ket{\psi}$, the expectation value of $H$ satisfies:
    \begin{equation}\label{eq:variational-principle}
        \bra{\psi} H \ket{\psi} \geq \lambda_0
    \end{equation}
    Where $\lambda_0$ is the minimum eigenvalue of $H$. Equality holds if and only if $\ket{\psi}$ is an eigenvector corresponding to $\lambda_0$.
\end{theorem}

\noindent
Thus, the ground-state energy $\lambda_0$ is the \textbf{lowest possible energy}, nothing can go below it. Obviously, this means that the ground state satisfies:
\begin{equation*}
    \bra{\psi_{\min}} H \ket{\psi_{\min}} = \underbrace{\lambda_{\min}}_{\lambda_0}
\end{equation*}

\highspace
\textcolor{Green3}{\faIcon{question-circle} \textbf{Why is the Variational Principle incredibly useful?}} Because it completely changes the problem. Originally we wanted to diagonalize the Hamiltonian to find its minimum eigenvalue, which is \textbf{computationally expensive}. Now we only need to minimize $\bra{\psi} H \ket{\psi}$ over all possible normalized states $\ket{\psi}$, which is \textbf{computationally feasible}. We shifted from a problem that is \textbf{exponentially hard} to a problem that is \textbf{polynomially hard}. Now, we need to solve an \textbf{optimization problem} that is still hard, but at least it is feasible.

\highspace
\begin{flushleft}
    \textcolor{Green3}{\faIcon{link} \textbf{Connection with the Ansatz}}
\end{flushleft}
The Ansatz is nothing more than a state generator. Instead of trying every possible quantum state (which is impossible), we define a family of states $\ket{\psi(\theta)}$ parameterized by $\theta$. The ansatz is literally a circuit that prepares these states.

\highspace
In other words, since the theorem needs a normalized state $\ket{\psi}$, \textbf{\emph{where does this state come from?}} It comes from the Ansatz. The Ansatz is a circuit that prepares a quantum state $\ket{\psi(\theta)}$ parameterized by $\theta$, then we measure the expectation value:
\begin{equation*}
    \bra{\psi(\theta)} H \ket{\psi(\theta)}
\end{equation*}
Simply put, the ansatz is a parameterized quantum circuit that generates trial quantum states. Thanks to the Variational Principle, VQE only needs to search among these trial states for the one with the smallest expected energy, instead of diagonalizing the exponentially large Hamiltonian.

\highspace
In summary, the Variational Quantum Eigensolver (VQE) is based on the\break Rayleigh-Ritz Variational Principle, which states that for any normalized quantum state $\ket{\psi}$, the expectation value of the Hamiltonian satisfies $\bra{\psi} H \ket{\psi} \ge \lambda_{\min}$, where $\lambda_{\min}$ is the minimum eigenvalue of the Hamiltonian. Equality is achieved only when $\ket{\psi}$ is the ground state. Therefore, if the ansatz can represent the ground state, minimizing the expectation value is equivalent to finding the ground-state energy of the molecule.